\section{Conclusion}
\label{sec:conclusion}
We show that the tension between robust concept erasure and utility preservation in diffusion models is a structural consequence of representation geometry, not merely a limitation of existing methods.
\Cref{thm:tradeoff} establishes a $\kappa$-scaled lower bound on collateral damage, and our experiments show that existing methods occupy different points on the resulting robustness--retention frontier: robust fine-tuning methods approach strong erasure but damage overlapping concepts, while inference-time methods preserve utility but leave residual leakage.

Our results suggest that this trade-off emerges naturally from overlapping representations. This indicates that the goal should not be to optimize erasure or retention in isolation, but rather to design $\kappa$-aware unlearning methods that apply displacement only along the activation directions necessary for erasure, instead of broadly perturbing shared adversarial subspaces. Our results also suggest that $\hat\kappa$ can serve as a deployment-time diagnostic: estimated from base-model activations alone, it predicts which targets admit clean erasure and which require a sharper Pareto trade-off.
%$\hat\kappa$ can serve as a deployment-time diagnostic: because it can be estimated from base-model activations, it helps predict which target concepts are amenable to clean erasure and which require accepting a sharper Pareto trade-off.

% \noindent{\bf The trade-off is unavoidable, not unsolvable.}
% \Cref{thm:tradeoff} is a lower bound, not a wall. Current methods cluster into three regimes (robust fine-tuning, standard fine-tuning, inference-time) corresponding to different effective displacement magnitudes $\Delta$, and that within each regime, performance varies in how closely the bound is approached. The  promising direction is therefore not "better erasure" or "better retention" in isolation, but \emph{$\kappa$-aware} displacement: applying the displacement only to the activation directions where it is required for erasure, rather than broadly across the spanned adversarial subspace as STEREO does.

% \noindent{\bf $\hat\kappa$ as a deployment-time diagnostic.}
% $\hat\kappa$ is estimated from the base model's activations on the target and a candidate neighbor pool, in a single forward pass per prompt. This estimate predicts not only the absolute severity of the trade-off but the qualitative regime (low-$\hat\kappa$ concepts admit operating points that high-$\hat\kappa$ concepts do not). When multiple concept descriptions are equivalent for a deployment goal (e.g., erasing "explicit content" via several candidate target concepts), $\hat\kappa$ provides a principled basis for choosing the version most amenable to clean erasure.

% \noindent{\bf Beyond diffusion.}
% The argument depends only on three properties: (i) the unlearning operation produces an activation displacement at the target concept, (ii) the activation map from prompts to representations is approximately Lipschitz, and (iii) target and neighbor concepts have geometrically overlapping activation regions. None are specific to diffusion, so the same $\kappa$-scaled trade-off should also apply to concept unlearning in autoregressive language models (where activations are tokens-in-context), classifiers (with penultimate features as activations), and multimodal models. We view our diffusion analysis as one instantiation of a more general phenomenon.

% Empirically, across seven unlearning methods and five target concepts spanning $\hat\kappa \in [0.27, 0.89]$, we observe (i) no method escapes the trade-off, (ii) trade-off severity scales monotonically with $\hat\kappa$, and (iii) operating regions stratify by mechanism type, with the most aggressive methods (STEREO, AdvUnlearn) approaching the predicted lower bound on high-$\hat\kappa$ concepts. 
% This perspective shifts concept unlearning from trying to bypass the trade-off to designing methods that approach the theoretical limit efficiently, while suggesting $\hat\kappa$ as a practical diagnostic for selecting targets and methods under deployment constraints.